\documentclass[11pt]{article}
\usepackage[margin=1in]{geometry}
\usepackage{amsmath,amssymb}
\usepackage[hidelinks]{hyperref}
\emergencystretch=3em

\newcommand{\Gammaq}{\Gamma_q}
\newcommand{\HR}{H_{\mathbb R}}

\title{Literature Status for the Two-Dimensional q-Araki-Woods Factoriality Gap in arXiv:1703.01317}
\author{Run fa\_banach\_001, lane 14}
\date{July 18, 2026}

\begin{document}
\maketitle

\section*{Status}

This note records that the factoriality and non-injectivity gap for
two-dimensional q-Araki-Woods algebras flagged in Avsec--Brannan--Wasilewski,
arXiv:1703.01317, is answered by Kumar--Skalski--Wasilewski,
arXiv:2301.08619.

No new theorem is claimed here. The purpose of this packet is to clear an
already-answered open-problem signal from the run queue.

\section*{Source Question}

In the introduction of arXiv:1703.01317, page 2, the authors survey known
structural results for q-Araki-Woods algebras. They state that for the
factoriality problem and non-injectivity, the remaining case is the family of
q-Araki-Woods algebras built from a two-dimensional Hilbert space. This is the
open-problem signal extracted from source line 140 of the local TeX source.

\section*{Later Answer}

The abstract of arXiv:2301.08619 says that the paper establishes
factoriality of q-Araki-Woods von Neumann algebras with at least two
generators in full generality, and also establishes non-injectivity and the
type of the factors.

The main theorem on page 2 states that if $q\in(-1,1)$, $\dim \HR\geq 2$,
and $(U_t)_{t\in\mathbb R}$ is a group of orthogonal transformations of
$\HR$, then $\Gammaq(\HR,U_t)$ is a non-injective factor, with type determined
by the closed subgroup of $\mathbb R_*^\times$ generated by the spectrum of
the generator of $(U_t)_{t\in\mathbb R}$. Theorem 2.5 on page 9 records
non-injectivity for $\dim \HR\geq 2$.

Therefore the two-dimensional case singled out in arXiv:1703.01317 is already
covered by later literature. The answer paper explicitly frames itself as a
full solution of the q-Araki-Woods factoriality question, so this packet is
classified as \texttt{literature\_already\_answered}.

\section*{Scope}

Resolved scope: factoriality and non-injectivity for q-Araki-Woods algebras
with underlying real Hilbert-space dimension at least two, including the
two-dimensional case named in arXiv:1703.01317.

Not resolved by this packet: the neighboring radial-multiplier conjecture in
arXiv:1703.01317, namely type-invariance and the expected isometric and
bijective form of the transference principle for radial multipliers.

\section*{Search Evidence}

The lane checked the run indexes for arXiv:1703.01317 and core keywords
including q-Araki-Woods, two-dimensional Hilbert space, factoriality,
non-injectivity, and radial multipliers. No prior packet for arXiv:1703.01317
was found. External arXiv/web checks identified arXiv:2301.08619 as the
decisive later answer. The later mixed q-Araki-Woods paper arXiv:2304.11108
also records that the complete factoriality and non-injectivity question for
the non-mixed q-Araki-Woods algebras had been settled by
Kumar--Skalski--Wasilewski.

\section*{References}

\begin{enumerate}
\item S. Avsec, M. Brannan, and M. Wasilewski,
\emph{Complete metric approximation property for q-Araki-Woods algebras},
arXiv:1703.01317; J. Funct. Anal. 274 (2018), no. 2, 544--572.
\item M. Kumar, A. Skalski, and M. Wasilewski,
\emph{Full solution of the factoriality question for q-Araki-Woods von
Neumann algebras via conjugate variables}, arXiv:2301.08619; Comm. Math.
Phys. 402 (2023), 157--167.
\item M. Kumar,
\emph{Conjugate variables approach to mixed q-Araki-Woods algebras:
Factoriality and non-injectivity}, arXiv:2304.11108; J. Math. Phys. 64
(2023), 093506.
\end{enumerate}

\end{document}
